\documentclass[11pt]{article}
\usepackage[a4paper,margin=22mm]{geometry}
\usepackage{amsmath,amssymb,amsthm,mathtools,mathrsfs}
\usepackage{graphicx}
\usepackage{microtype}
\usepackage{enumitem}
\usepackage{xcolor}
\usepackage[colorlinks=true,linkcolor=blue!55!black,citecolor=blue!55!black,urlcolor=blue!55!black]{hyperref}

\newtheorem{theorem}{Theorem}
\newtheorem{lemma}[theorem]{Lemma}
\newtheorem{corollary}[theorem]{Corollary}
\newcommand{\T}{\mathbb T}
\newcommand{\D}{\mathbb D}
\newcommand{\C}{\mathbb C}
\newcommand{\Mz}{\mathscr M_z}
\newcommand{\Res}{\operatorname{Res}}
\newcommand{\Disc}{\operatorname{Disc}}

\title{Generic Rank-Three Failure of Cauchy-Dual Subnormality\\
\large A full-dimensional partial answer to the finite-support characterization problem}
\author{Candidate partial result for arXiv:2309.03588}
\date{26 August 2026}

\begin{document}
\maketitle

\begin{center}
\fbox{\parbox{0.92\linewidth}{\textbf{Status: candidate substantial partial result, likely valid.}
For measures supported at exactly three points, the packet proves that the
subnormal locus is contained in a Lebesgue-null, nowhere-dense real-analytic
exceptional set.  It also proves non-subnormality throughout a full
neighborhood of every equal-mass equilateral measure.  The complete
classification on the exceptional set remains open.}}
\end{center}

\section{Source problem and claimed contribution}

Khasnis and Sholapurkar state the finite-support Cauchy-dual problem in
Section~1.2, Problem~1, on source PDF p.~5 \cite{KS}:

\begin{center}
\includegraphics[width=0.98\linewidth]{figures/open_problem_crop.png}
\end{center}

\noindent\textbf{Exact transcription.}
``Characterize finitely supported, positive, Borel measures $\mu$ on unit
circle $\T$ such that the Cauchy dual of $M_z$ on $D(\mu)$ is subnormal.''

For ordered three-point measures, normalize rotation and cyclic order by
writing
\begin{equation}\label{eq:chart}
 \mu_\lambda=c_1\delta_1+c_2\delta_{e^{i\theta_2}}
       +c_3\delta_{e^{i\theta_3}},\qquad
 \lambda\in\Omega,
\end{equation}
where
\[
 \Omega=\{(\theta_2,\theta_3,c_1,c_2,c_3):
       0<\theta_2<\theta_3<2\pi,\ c_j>0\}.
\]
This connected open set is a full parameter chart for rank-three measures,
up to rotation and relabeling.

\begin{theorem}[Generic rank-three non-subnormality]\label{thm:generic}
There are two nonzero real-analytic functions $\mathcal D,\mathcal R$ on
$\Omega$ such that, if the Cauchy dual of $\Mz$ on $D(\mu_\lambda)$ is
subnormal, then
\[
 \mathcal D(\lambda)\mathcal R(\lambda)=0.
\]
Consequently the subnormal rank-three measures form a Lebesgue-null and
nowhere-dense subset of every chart \eqref{eq:chart}.  Equivalently, an open
dense set of rank-three measures, and almost every such measure in the
natural angle--mass coordinates, has a non-subnormal Cauchy dual.
\end{theorem}

\begin{corollary}[Robust equilateral obstruction]\label{cor:stability}
Fix $c>0$, $\xi\in\T$, and $\omega=e^{2\pi i/3}$.  There is an
$\varepsilon>0$ such that, whenever the atoms can be labeled so that
\[
 |\zeta_j-\xi\omega^j|<\varepsilon,
 \qquad |c_j-c|<\varepsilon \quad (j=0,1,2),
\]
the Cauchy dual on
$D(\sum_{j=0}^2c_j\delta_{\zeta_j})$ is not subnormal.
Thus the known equilateral example lies inside a full-dimensional open set
of arbitrary support triples and arbitrary positive masses.
\end{corollary}

\section{Rational model and the forbidden-atom criterion}

For a finite positive measure $\mu$ on $\T$, $D(\mu)$ denotes the
corresponding Dirichlet-type space, and $\Mz f=zf$ is a cyclic analytic
$2$-isometry.  Its Cauchy dual is
$\Mz'=\Mz(\Mz^*\Mz)^{-1}$.

The rational de Branges--Rovnyak model of Chavan--Ghara--Reza and Costara
\cite{CGR,Costara} gives, for every $\lambda\in\Omega$, a monic outer
polynomial
\[
 q_\lambda(z)=\prod_{r=1}^3(z-\alpha_r),\qquad |\alpha_r|>1,
\]
and a row Schur function
\[
 B=(p_1/q_\lambda,p_2/q_\lambda,p_3/q_\lambda)
\]
such that $D(\mu_\lambda)=\mathcal H(B)$ isometrically.  Put
\begin{equation}\label{eq:G}
 G_\lambda(z,w)=\sum_{j=1}^3p_j(z)\overline{p_j(w)}.
\end{equation}
The polynomial $G_\lambda$ is intrinsic even though a matrix factor producing
the individual $p_j$ is not unique.

When the roots of $q_\lambda$ are simple, set
\begin{equation}\label{eq:kappa}
 \kappa_{rt}=
 \frac{G_\lambda(\alpha_r,\alpha_t)}
 {\alpha_r^2\overline{\alpha_t}^{\,2}
  q_\lambda'(\alpha_r)\overline{q_\lambda'(\alpha_t)}}.
\end{equation}
Corollary~4.1 and Remark~4.2 of \cite{CGR} give the following necessary
condition: if $\Mz'$ is subnormal, then
\begin{equation}\label{eq:nu}
 \nu_\lambda=\sum_{r,t=1}^3\kappa_{rt}
 \delta_{1/(\alpha_r\overline{\alpha_t})}
\end{equation}
is a positive measure supported on $[0,1]$.  Equal atom locations in
\eqref{eq:nu} must be grouped before their coefficients are tested.  In
particular, every grouped coefficient outside $[0,1]$ must vanish.

The next observation sharpens Corollaries~2.3 and 4.3 of \cite{CGR} by
removing every hypothesis on the pole products.

\begin{lemma}[Simple-root orthogonality criterion]\label{lem:orthogonality}
If the roots $\alpha_1,\alpha_2,\alpha_3$ are distinct, then the Cauchy dual
is subnormal if and only if
\begin{equation}\label{eq:crosszero}
 G_\lambda(\alpha_r,\alpha_t)=0\qquad(r\ne t).
\end{equation}
\end{lemma}

\begin{proof}
Put $u_r=1/\alpha_r$ and let $c_{rt}$ denote the quantity in
\eqref{eq:kappa}.  Theorem~2.1 of \cite{CGR} says that subnormality makes,
for every integer $\ell\ge1$, the infinite matrix
\[
 \left(\sum_{r,t=1}^3 c_{rt}
 (1-u_r\overline{u_t})^\ell u_r^m\overline{u_t}^{\,n}
 \right)_{m,n\ge0}
\]
formally positive semidefinite.  Compress its leading $3\times3$ block and
conjugate by the inverse of the Vandermonde matrix
$(u_r^m)_{0\le m\le2,,1\le r\le3}$.  This gives
\[
 C_\ell=\big(c_{rt}(1-u_r\overline{u_t})^\ell\big)_{r,t=1}^3\succeq0.
\]
If $G(\alpha_r,\alpha_t)\ne0$, positivity of the numerator Gram kernel makes
both corresponding diagonal values positive.  The $(r,t)$ principal minor
of $C_\ell$ would therefore require
\[
 c_{rr}c_{tt}\big((1-|u_r|^2)(1-|u_t|^2)\big)^\ell
 \ge |c_{rt}|^2|1-u_r\overline{u_t}|^{2\ell}.
\]
But
\[
 |1-u_r\overline{u_t}|^2
 -(1-|u_r|^2)(1-|u_t|^2)=|u_r-u_t|^2>0,
\]
so the displayed inequality fails for all sufficiently large $\ell$.
Thus every cross value vanishes.  Conversely, Corollary~2.3 of \cite{CGR}
proves subnormality when all cross values vanish.
\end{proof}

\subsection*{Proof Intuition}

There are two independent algebraic obstructions.  The discriminant of the
outer factor detects repeated roots.  Away from that hypersurface,
Lemma~\ref{lem:orthogonality} forces all cross-Gram values to vanish, and one
iterated resultant is the product of all nine Gram values
$G(\alpha_r,\alpha_t)$.  Neither analytic function is identically zero: the
exact equilateral calculation witnesses both facts.  Proper real-analytic
zero sets have measure zero and empty interior.

The neighborhood result uses more information at the equilateral point.  Its
six off-diagonal residues form two clusters, and all three residues in either
cluster have the same nonzero value.  Under a small perturbation, even if
atoms collide, no nonempty sum of nearby residues can cancel.  Hence a
forbidden atom remains outside $[0,1]$.

\section{Analytic dependence and the two obstructions}

\begin{lemma}[Analytic rational data]\label{lem:analytic}
The coefficients of $q_\lambda$ and $G_\lambda$ depend real-analytically on
$\lambda\in\Omega$.
\end{lemma}

\begin{proof}
Let $s_\lambda(z)=\prod_{j=1}^3(z-\zeta_j)$.  The outer factor is uniquely
determined, together with $\gamma_\lambda>0$, by
\begin{equation}\label{eq:fejer}
 |s_\lambda(z)|^2+
 \sum_{j=1}^3c_j\left|\frac{s_\lambda(z)}{z-\zeta_j}\right|^2
 =\gamma_\lambda|q_\lambda(z)|^2,\qquad z\in\T.
\end{equation}
The Laurent coefficients on the left are real-analytic in $\lambda$.
The coefficient map $(\gamma,q)\mapsto\gamma|q|^2|_\T$, with $q$ monic
and zero-free on $\overline\D$, has invertible linearization.  Indeed, a
kernel vector would give on $\T$
\[
 a+2\operatorname{Re}(h/q)=0,
\]
where $a\in\mathbb R$ and $\deg h<\deg q$.  The analytic function $h/q$ on
$\D$ has constant real part, hence is constant; the degree restriction then
forces $h=0$, and also $a=0$.  The real-analytic implicit-function theorem
and uniqueness of the outer factor give analytic dependence of $q$.

Costara's kernel formula expresses $G$ using $q$, $s$, the rational functions
\[
 O=\gamma^{-1/2}e^{i\vartheta}s/q,\qquad
 f_j=\frac{O}{O'(\zeta_j)(z-\zeta_j)},
\]
and the inverse of the positive Gram matrix
$(\langle f_i,f_j\rangle)_{i,j=1}^3$; see Theorem~6.1 and Remark~6.2 of
\cite{CGR}.  Every denominator stays nonzero on $\Omega$, and the phase
normalization cancels in the intrinsic product \eqref{eq:G}.  These finite
algebraic operations and a matrix inverse preserve real-analytic dependence.
\end{proof}

Define the outer-root discriminant
\begin{equation}\label{eq:D}
 \mathcal D(\lambda)=\Disc_z q_\lambda(z).
\end{equation}
Thus $\mathcal D\ne0$ exactly when the three outer roots are distinct.

Write $\widehat G_\lambda(z,y)$ for the polynomial obtained from
$G_\lambda(z,w)$ by replacing $\overline w$ by $y$, and put
$q_\lambda^c(y)=\overline{q_\lambda(\overline y)}$.  Define
\begin{equation}\label{eq:R}
 \mathcal R(\lambda)=
 \Res_z\!\left(q_\lambda(z),
   \Res_y\!\left(q_\lambda^c(y),\widehat G_\lambda(z,y)\right)
 \right).
\end{equation}
Both polynomials $q_\lambda$ and $q_\lambda^c$ are monic, so the elementary
root formula for resultants gives
\begin{equation}\label{eq:Rproduct}
 \mathcal R(\lambda)=
 \prod_{r,t=1}^3G_\lambda(\alpha_r,\alpha_t)
 =\prod_{r=1}^3G_\lambda(\alpha_r,\alpha_r)
  \prod_{r<t}|G_\lambda(\alpha_r,\alpha_t)|^2.
\end{equation}
In particular $\mathcal R$ is real-valued.  Lemma~\ref{lem:analytic} shows
that $\mathcal D$ and $\mathcal R$ are real-analytic on $\Omega$.

\section{The two nonzero witnesses}

\begin{lemma}[The resultant is nontrivial]\label{lem:Rnonzero}
$\mathcal R$ is not identically zero on $\Omega$.
\end{lemma}

\begin{proof}
Fix $c>0$ and take the equal-mass equilateral measure
$c(\delta_1+\delta_\omega+\delta_{\omega^2})$.  Root-of-unity
orthogonality in \eqref{eq:fejer} gives
\[
 q(z)=z^3-\rho,\qquad
 \rho+\rho^{-1}=2+9c,\qquad \rho>1.
\]
The circulant Costara Gram matrix diagonalizes in the Fourier basis and yields
\begin{equation}\label{eq:P}
 G(z,w)=P(z\overline w),\qquad
 P(x)=\sum_{k=1}^3A_kx^k,\qquad
 A_k=\frac{3\rho(\rho-1)^2}
 {[3+(k-1)(\rho-1)][3+k(\rho-1)]}>0.
\end{equation}
This is the $N=3$ specialization of the same rational-kernel calculation
used in \cite{JPS}; it also follows directly from Theorem~6.1 of \cite{CGR}.

Put $a=\rho^{1/3}$.  The roots of $q$ are $a,a\omega,a\omega^2$.
The diagonal value $P(a^2)$ is positive.  If $P(a^2\omega)=0$, real
coefficients also give $P(a^2\omega^2)=0$.  Hence $P(x)/x$ would be a
multiple of
\[
 (x-a^2\omega)(x-a^2\omega^2)=x^2+a^2x+a^4,
\]
forcing $A_2/A_1=a^{-2}=\rho^{-2/3}$.  But \eqref{eq:P} gives
\[
 \frac{A_2}{A_1}=\frac{3}{3+2(\rho-1)}<\rho^{-2/3}.
\]
For the strict inequality, write $a>1$ and note that
\[
 1+\frac{2(a^3-1)}3-a^2
\]
vanishes at $a=1$ and has derivative $2a(a-1)>0$.  Thus both off-diagonal
values are nonzero.  Every factor in \eqref{eq:Rproduct} is nonzero, so
$\mathcal R\ne0$ at this equilateral point.
\end{proof}

\begin{lemma}[The outer-root discriminant is nontrivial]\label{lem:Dnonzero}
$\mathcal D$ is not identically zero on $\Omega$.
\end{lemma}

\begin{proof}
At the equal-mass equilateral point used above,
$q(z)=z^3-\rho$ with $\rho>1$.  Its three roots are distinct, so
$\mathcal D\ne0$ there.
\end{proof}

\section{Proofs of the main results}

\begin{proof}[Proof of Theorem~\ref{thm:generic}]
Suppose $\mathcal D(\lambda)\ne0$.  Then the roots of $q_\lambda$ are
distinct.  If the Cauchy dual were subnormal,
Lemma~\ref{lem:orthogonality} would force \eqref{eq:crosszero}; then
\eqref{eq:Rproduct} would give $\mathcal R(\lambda)=0$.  Therefore every
subnormal parameter belongs to
\[
 \{\mathcal D=0\}\cup\{\mathcal R=0\}.
\]
By Lemmas~\ref{lem:Rnonzero} and \ref{lem:Dnonzero}, these are the zero sets
of two nonzero real-analytic functions on the connected open set $\Omega$.
A proper real-analytic zero set has Lebesgue measure zero and empty interior
\cite{KP}.  Their finite union has the same properties, proving the theorem.
\end{proof}

\begin{proof}[Proof of Corollary~\ref{cor:stability}]
Rotation is harmless, so take $\xi=1$.  At the equal-mass equilateral point,
label $\alpha_r=a\omega^r$, $r=0,1,2$.  If $d=r-t\not\equiv0\pmod3$, then
\eqref{eq:kappa} and $q'(\alpha_r)=3\alpha_r^2$ give the common residue
\begin{equation}\label{eq:base-residue}
 \kappa_{rt}=K_d:=
 \frac{P(a^2\omega^d)}{9\rho^2a^2\omega^d}\ne0.
\end{equation}
Thus the three ordered pairs with difference $d=1$ have one common nonzero
residue, and the three with difference $d=2$ have another.  Their atom
locations are the two nonreal points
$a^{-2}\omega^{-1}$ and $a^{-2}\omega^{-2}$.

The outer roots stay simple and can be labeled continuously under a small
perturbation.  Lemma~\ref{lem:analytic} and \eqref{eq:kappa} show that the
six residues and their atom locations vary continuously.  Choose disjoint
neighborhoods $U_1,U_2$ of the two base locations, both disjoint from
$[0,1]$.  Shrink the parameter neighborhood until every residue in cluster
$d$ differs from $K_d$ by less than $|K_d|/2$.

Any collision group inside $U_d$ is a nonempty sum of, say, $m$ such
residues.  Its distance from $mK_d$ is less than $m|K_d|/2$, so the grouped
coefficient is nonzero.  Hence the measure \eqref{eq:nu} has a nonzero atom
outside $[0,1]$, contradicting the necessary condition for subnormality.
This holds for every sufficiently small perturbation, including parameters
where some pole products collide.
\end{proof}

\section{Verification, limitations, and novelty audit}

The accompanying script \texttt{code/verify\_rank\_three.py} independently
constructs the degree-six self-inversive polynomial, selects the three outer
roots, reconstructs Costara's Gram matrix and the intrinsic polynomial $G$,
and evaluates the nine pole products, six cross values, and six forbidden
residues.  It passes three cases: the equal-mass equilateral model, a fully
nonsymmetric perturbation, and a small-mass generic witness.  In the latter
two cases all nine pole products are separated and every cross value is
nonzero.  These floating-point checks are sanity tests, not proof.

\textbf{Adversarial verification focus.}
The most delicate points are (i) the real-analytic dependence in
Lemma~\ref{lem:analytic} and (ii) the Vandermonde compression and
large-$\ell$ principal-minor argument in Lemma~\ref{lem:orthogonality}.
The local stability proof separately handles pole-product collisions by
grouping the residues exactly as required by Remark~4.2 of \cite{CGR}.

\textbf{Scope limitation.}
The theorem does not classify the proper real-analytic exceptional set
$\{\mathcal D\mathcal R=0\}$.  Subnormal points, if any beyond the known boundary
and lower-rank cases, must lie there, but membership in that set is not
sufficient for subnormality.  The full arbitrary-three-point conjecture and
the characterization for four or more atoms remain open.

\textbf{Bounded novelty audit.}
The run indexes and primary-source arXiv searches through 26 August 2026 used
arXiv:2103.10059, 2309.03588, 2510.14004, and 2511.04565; the authors Khasnis,
Phatak, Sholapurkar, Chavan, Ghara, and Reza; and the phrases ``Cauchy dual
subnormality,'' ``rank three,'' ``three distinct points,'' ``almost every,''
``generic,'' and ``perturbation.''  The search found the rational-symbol
criterion, the original rank-two example, the later rank-two classification,
and the equal-mass equilateral rank-three example \cite{CGR,KS,JPS}; it found
no generic, full-measure, or perturbative rank-three theorem.  Novelty is
plausible, not certified.

\textbf{Human-review recommendation.}
Send to an operator theorist familiar with rational de Branges--Rovnyak
spaces.  If the simple-root criterion and analytic-dependence lemma pass
review, the generic
theorem is a checkable substantial partial answer; the exceptional
set is then the natural target for a full classification.

\begin{thebibliography}{9}
\bibitem{KS}
M.~N. Khasnis and V.~M. Sholapurkar,
\emph{Cyclic analytic $2$-isometry of finite rank and Cauchy dual
subnormality problem}, arXiv:2309.03588v2 (2024).

\bibitem{CGR}
S. Chavan, S. Ghara, and M.~R. Reza,
\emph{The Cauchy dual subnormality problem via de Branges--Rovnyak spaces},
Studia Math. \textbf{265}(3) (2022), 315--351; arXiv:2103.10059.

\bibitem{Costara}
C. Costara,
\emph{Reproducing kernels for Dirichlet spaces associated to finitely
supported measures}, Complex Anal. Oper. Theory \textbf{10} (2016),
1277--1293.

\bibitem{JPS}
S.~A. Joshi, G.~M. Phatak, and V.~M. Sholapurkar,
\emph{An example of a cyclic analytic $2$-isometry with defect operator of
rank $3$, whose Cauchy dual is not subnormal}, arXiv:2511.04565 (2025).

\bibitem{KP}
S.~G. Krantz and H.~R. Parks,
\emph{A Primer of Real Analytic Functions}, second ed., Birkh\"auser, 2002.
\end{thebibliography}

\end{document}
